\documentclass[12pt,a4paper]{article}
\usepackage[utf8]{inputenc}
\usepackage[vietnamese]{babel}
\usepackage[T5]{fontenc}
\usepackage{amsmath,amssymb,amsfonts}
\usepackage{fancyhdr}
\usepackage{geometry}
\usepackage{tcolorbox}

\geometry{
	a4paper,
	left=20mm,
	right=20mm,
	top=20mm,
	bottom=20mm
}

\pagestyle{fancy}
\fancyhf{}
\lhead{\small\textit{Bổ đề Hình học: Định lý Thales \& Phương tích}}
\rhead{\small\textbf{Integra Typesetter}}
\rfoot{\small Trang \thepage}
\renewcommand{\headrulewidth}{0.4pt}
\renewcommand{\footrulewidth}{0.4pt}

\begin{document}

\begin{center}
	{\Large \textbf{CÁC BỔ ĐỀ HÌNH HỌC CỐT LÕI}}\\[2mm]
	{\small \textit{(Định lý Thales, Tam giác đồng dạng và Phương tích đường tròn)}}
\end{center}
\vspace{4mm}

\begin{tcolorbox}[colback=blue!5!white,colframe=blue!75!black,title=\textbf{Bổ đề 1: Định lý Thales và Tam giác đồng dạng với các đường song song}]
Cho tam giác $TNC$ có $M \in TN, D \in TC$. 
\begin{enumerate}
    \item Nếu $MD \parallel NC$ thì $\dfrac{TD}{TC} = \dfrac{TM}{TN} = \dfrac{MD}{NC}$.
    \item Nếu $DN \parallel MB$ (với $B \in TD$) thì $\dfrac{TB}{TD} = \dfrac{TM}{TN} = \dfrac{BM}{DN}$.
\end{enumerate}
\end{tcolorbox}

\vspace{3mm}

\begin{tcolorbox}[colback=teal!5!white,colframe=teal!75!black,title=\textbf{Bổ đề 2: Phương tích và Tiêu chuẩn tiếp xúc của đường tròn}]
Cho ba điểm $E, D, F$ không thẳng hàng và điểm $T$ thuộc đường thẳng $EF$ ($T$ nằm ngoài đoạn $EF$). Đường thẳng qua $T$ và $D$ tiếp xúc với đường tròn ngoại tiếp tam giác $DEF$ tại $D$ khi và chỉ khi:
\[
TD^2 = TE \cdot TF \iff \widehat{TDE} = \widehat{TFD} \iff \Delta TDE \sim \Delta TFD \text{ (c-g-c)}
\]
\end{tcolorbox}

\end{document}
